% 组合数学（综述）
% license CCBYSA3
% type Wiki

本文根据 CC-BY-SA 协议转载翻译自维基百科\href{https://en.wikipedia.org/wiki/Combinatorics}{相关文章}。

组合数学是数学的一个分支，主要研究计数（counting）——既作为获得结果的方法，又作为研究有限结构某些性质的目的。它与数学的许多其他领域密切相关，并在从逻辑到统计物理、从进化生物学到计算机科学等众多领域中有广泛的应用。

组合数学以其所处理问题的广泛性而著称。组合问题出现在纯数学的许多分支中，尤其是在代数学、概率论、拓扑学与几何学中\(^\text{[1]}\) ，同时也出现在众多应用领域。  
历史上，许多组合问题往往是独立提出的，作为在某个数学背景下出现的特定问题的特设性解法（ad hoc solutions）。然而，在二十世纪后期，强大而一般的理论方法被发展出来，使得组合数学成为一门独立的数学分支\(^\text{[2]}\)。组合数学中最古老且最易于理解的部分之一是图论（graph theory），它本身就与其他数学领域有着大量自然的联系。在计算机科学中，组合数学被广泛用于推导算法分析中的公式与估计。

